\documentclass[10pt,a4paper]{article}
\usepackage[margin=1in]{geometry}
\usepackage[T1]{fontenc}
\usepackage[utf8]{inputenc}
\usepackage{lmodern}
\usepackage{enumitem}
\setlist[itemize]{itemsep=1pt,topsep=2pt}
\setlength{\parindent}{0pt}
\setlength{\parskip}{4pt}
\newcommand{\verticalflowfour}[4]{\begin{center}\small\fbox{\strut #1}\\$\downarrow$\\\fbox{\strut #2}\\$\downarrow$\\\fbox{\strut #3}\\$\downarrow$\\\fbox{\strut #4}\end{center}}
\newcommand{\branchdiagram}{\begin{center}\small\fbox{\strut New lead}\ $\rightarrow$\ \fbox{\strut Rules + OpenClaw}\ $\rightarrow$\ \begin{tabular}{c}\fbox{\strut Reply draft}\\[2pt]\fbox{\strut Human review}\end{tabular}\end{center}}
\newcommand{\metricdiagram}{\begin{center}\small New enquiry \quad$\bullet\longrightarrow\bullet\longrightarrow\bullet$\quad Checked reply\\[-2pt]Target: under 5 min\end{center}}
\newcommand{\phasediagram}{\begin{center}\small\fbox{\strut Build MVP}\ $\rightarrow$\ \fbox{\strut Test and improve}\ $\rightarrow$\ \fbox{\strut Add later}\end{center}}

\title{OpenClaw-Based Real Estate Lead Desk for Small Property Teams}
\author{Aman Kapoor \quad Vidushi Jain}
\date{August 17, 2026}

\begin{document}
\maketitle

\section{Higher-order goal}
This project helps small real-estate teams respond to new customer enquiries faster while keeping the human agent in control.

\section{Time-to-value}
\begin{itemize}
    \item We will first build the most useful flow: add an enquiry, collect the main details, find matching properties, and prepare a reply.
    \item We will keep the first version small and free. It will not include paid WhatsApp APIs, CRM connections, or website scraping.
    \item We will test whether an agent can reply faster with our system than by reading chats and searching a spreadsheet manually.
\end{itemize}

\section{Problem Statement}
Small real-estate teams receive messages from calls, WhatsApp, forms, and social media. Important details such as budget, location, and property type are spread across these messages, so agents must search manually before replying.

\section{Proposed Solution}
\subsection{Overview}
We will build a small web app called a \textbf{real-estate lead desk}. A staff member pastes a customer enquiry into the app, and the final reply is always checked by a human agent.
\verticalflowfour{Lead inbox}{OpenClaw reads needs}{Property matches}{Human reply}

\subsection{Operational constraints for an educational, tier-2/3 setting}
\begin{itemize}
    \item The app should work on normal student laptops and common phones.
    \item We will use simple matching rules instead of training a complex AI model.
    \item The first version will run locally with SQLite and CSV files, so it can be built for free.
\end{itemize}

\section{Solution Approach}
This is a focused student project, not a full CRM system. It supports one job: helping an agent prepare a reply to a new lead.
\branchdiagram

\subsection{Duplicate/quality assistance}
The app will check whether a new enquiry looks similar to an older one. It can compare the contact detail, location, budget, and property type. It will also show missing details clearly. The system will only suggest possible duplicates; it will never remove or merge leads by itself. OpenClaw will not give legal, financial, or mortgage advice. Such questions will be shown to the human agent.

\paragraph{Free data strategy.}
For the demo, we will create a fictional property list with 20--30 entries. Each property will have simple fields such as locality, property type, bedrooms, budget range, and availability. We will also create 30--50 fictional customer enquiries. This lets us test the system without using real phone numbers or private messages. The Ames Housing dataset may be used only as a public example for property fields; it will not be shown as real local listings.

\subsection{Web architecture}
The frontend will show the lead inbox, property matches, and reply draft. The backend will save leads, read the CSV file, run the OpenClaw flow, and apply the matching rules. SQLite will store the test data. This is enough for a working student prototype and does not need paid hosting.

\subsection{CI/CD and fast delivery enabler}
We will use GitHub and GitHub Actions for version control and basic tests. The app will run on our laptops. If we have extra time, we can put the frontend on a free hosting service. We will not build paid messaging, SMS, user accounts, or full CRM features in the first version.

\section{Evaluation Criterion}
\subsection{Primary evaluation metric}
\textbf{Time-to-ready-reply (TRR)} measures the time from adding a customer enquiry to getting a checked reply draft. Our target is a median TRR of less than five minutes.
\metricdiagram

\subsection{Secondary evaluation metrics}
\begin{itemize}
    \item \textbf{Complete summary:} are the main fields filled in or clearly marked as missing?
    \item \textbf{Good matches:} do the suggested properties fit the location, budget, and type requested?
    \item \textbf{Useful reply:} can a reviewer edit the reply quickly and send it?
    \item \textbf{Safe behaviour:} does the app ask for human help when the question is unclear or needs professional advice?
    \item \textbf{Reliability:} does the full flow work from adding a lead to saving the reply draft?
\end{itemize}

\subsection{Pilot validation plan}
We will test the app using 30--50 fictional enquiries from buyers, tenants, and sellers. Two to five people will review some of the summaries, matches, and reply drafts. We will compare the app with the manual method of reading the message, searching a spreadsheet, and writing a reply.

\section{Scalability}
The first version uses separate lead, matching, workflow, and data modules. Later, its CSV file can move to a database or CRM connection; for now, local SQLite is enough.

\section{Engine availability heuristic}
\begin{itemize}
    \item OpenClaw for the assistant flow, with a simple template backup if it is unavailable.
    \item A lightweight web framework and backend API.
    \item SQLite and CSV files for the data.
    \item A fictional property list and fictional customer enquiries for testing.
    \item GitHub and GitHub Actions for code and tests.
\end{itemize}

\section{Project scope and deliverables}
\subsection{Initial deliverable}
\begin{itemize}
    \item Lead inbox where the user can paste or type a customer enquiry.
    \item CSV import for a fictional property list.
    \item OpenClaw detail extraction with a simple template backup.
    \item Property matching, lead summary, and editable reply draft.
    \item SQLite storage and basic automated tests.
\end{itemize}

\subsection{Subsequent deliverables}
\begin{itemize}
    \item Search and filters for locality, budget, property type, and follow-up status.
    \item Simple response templates for a specific agency.
    \item Copy-to-clipboard for WhatsApp or email.
    \item A short report, test results, and final demonstration.
\end{itemize}
\phasediagram

\section{Risks and mitigations}
\begin{itemize}
    \item \textbf{Wrong AI output:} show the original enquiry, mark missing fields, and let a human check every reply.
    \item \textbf{Project becomes too large:} do not add logins, scraping, CRM links, notifications, or analytics in the first version.
    \item \textbf{Data privacy:} use only fictional property data and fictional customer enquiries for the academic demo.
    \item \textbf{OpenClaw is unavailable:} use a rule-based reply template so the project still works.
    \item \textbf{Two-student workload:} Aman can focus on backend, data, and OpenClaw; Vidushi can focus on frontend, testing, and documentation. Both will test the full system together.
\end{itemize}

\section{Summary}
Aman Kapoor and Vidushi Jain will build a free, practical lead desk that helps agents prepare faster replies while keeping people in control.

\end{document}
